\documentclass[11pt]{article}
\usepackage[margin=1in]{geometry}
\usepackage{xurl}\usepackage{amsmath,amssymb,amsthm,booktabs,longtable,hyperref}
\newtheorem{theorem}{Theorem}[section]\newtheorem{proposition}[theorem]{Proposition}
\newtheorem{lemma}[theorem]{Lemma}\newtheorem{corollary}[theorem]{Corollary}
\theoremstyle{definition}\newtheorem{definition}[theorem]{Definition}\newtheorem{remark}[theorem]{Remark}
\newcommand{\PP}{\mathbb{P}}\newcommand{\CC}{\mathbb{C}}\newcommand{\QQ}{\mathbb{Q}}
\title{Lean 4 Jacobian-Ring Certificates for the Fermat Quintic Fourfold:\\ Smoothness, Hilbert Series $(1+t+t^2+t^3)^6$,\\ and Primitive Hodge Numbers $(0,120,580,120)$}
\author{Benjamin Stanley Frohman\\ \small ORCID 0009-0006-7068-3718}
\date{November 2026 (version 1.3.0 of the FermatQuinticJacobian record)}
\sloppy\begin{document}\maketitle
\begin{abstract}
For $G=x_0^5+\dots+x_5^5$ we give machine-checked (Lean~4, Mathlib) proofs that the partial derivatives
$\partial_iG=5x_i^4$ have no common zero away from the origin in characteristic $0$, so $X=V(G)\subset\PP^5$ is
smooth, and we compute in Lean the Hilbert series $(1+t+t^2+t^3)^6$ of the Jacobian ring
$R=\CC[x]/(x_0^4,\dots,x_5^4)$: 19 coefficients summing to 4096, with $R_4=120$, $R_9=580$, $R_{14}=120$.
Via Griffiths' residue theorem (cited) these are the primitive Hodge numbers $h^{3,1},h^{2,2}_{\rm prim},h^{1,3}$,
and $h^{4,0}=0$. We explain how the series is encoded as a decidable list computation, cross-check it against
the Euler characteristic and an independent Singular computation, and apply the same Jacobian-ring analysis
to the Dwork-type members $\psi=1$ (singular, Lean) and $\psi=2$ (smooth, Lean). All declarations use only the
standard axioms; no claim about algebraic cycles beyond cited results is made.
\end{abstract}
\tableofcontents
\section{Introduction}
The Jacobian ring is where the Hodge theory of a smooth hypersurface becomes finite algebra. For the Fermat
quintic fourfold this algebra is explicit, which makes it a clean target for formal verification. Earlier
versions of this record contained a two-page note certifying the partial-derivative formulas (1.0.1, 1.1.0) and
a short update note (1.2.0). This version is the full paper and supersedes them. The focus here is the
Jacobian ring itself: its structure, the formal computation of its Hilbert series, and exactly where Lean
stops and Griffiths' theorem begins.
\section{Setup and smoothness}\label{sec:setup}
Let $K$ be a field of characteristic $0$ (for Hodge theory, $K=\CC$) and
\[ G=\sum_{i=0}^5 x_i^5\in K[x_0,\dots,x_5],\qquad X=V(G)\subset\PP^5_K. \]
$G$ is homogeneous of degree $5$ (Lean: \texttt{QuinticHost.G\_isHomogeneous}), so $X$ is a hypersurface
of dimension 4. Euler's identity $5G=\sum_i x_i\,\partial_iG$ (Lean: \texttt{FermatQuintic.euler\_homogeneous})
shows that a common zero of the partials lies on $X$.

\begin{theorem}[Lean: \texttt{QuinticHost.pderiv\_G}, \texttt{QuinticHost.smooth}]
$\partial_iG=5x_i^4$, and if $\mathrm{char}\,K=0$ the six partials vanish simultaneously only at $x=0$.
Consequently $X$ is smooth (Jacobian criterion).
\end{theorem}
\begin{proof}
$5x_i^4=0$ with $5\neq0$ in $K$ gives $x_i^4=0$, hence $x_i=0$ in the field $K$.
\end{proof}
The Jacobian criterion itself (a projective hypersurface is smooth iff the partials have no common zero
off the origin) is standard and is not formalised; Lean proves the polynomial statement.

Basic invariants: $K_X=\mathcal O_X(5-6)=\mathcal O_X(-1)$, so $X$ is Fano of index 1. By Lefschetz,
$b_0=b_2=b_6=b_8=1$ and $b_1=b_3=b_5=b_7=0$. The topological Euler characteristic of a smooth degree-$d$
hypersurface in $\PP^5$ is $\chi=\frac{(1-d)^6-1}{d}+6$; for $d=5$ this is $\frac{4096-1}{5}+6=825$,
so $b_4=825-4=821$ and the primitive part has dimension $820$.

\section{The Jacobian ring and its Hilbert series}\label{sec:jac}
The Jacobian ideal is $J=(\partial_0G,\dots,\partial_5G)=(x_0^4,\dots,x_5^4)$ and the Jacobian ring is
\[ R=\CC[x_0,\dots,x_5]/J \;=\;\bigotimes_{i=0}^5 \CC[x_i]/(x_i^4). \]
Since $x_0^4,\dots,x_5^4$ is a regular sequence, $R$ is a graded Artinian complete intersection with monomial
basis $\{x^a: 0\le a_i\le3\}$, so
\[ H_R(t)=\Big(\frac{1-t^4}{1-t}\Big)^6=(1+t+t^2+t^3)^6,\qquad \dim R=4^6=4096, \]
with socle in degree $6\cdot 3=18$ and the symmetry $\dim R_k=\dim R_{18-k}$ (Macaulay duality).
For an arbitrary smooth quintic in six variables the partials still form a regular sequence of six forms
of degree 4, so the same Hilbert series holds (cited, not formalised).
\begin{table}[h]\centering\small
\begin{tabular}{l*{10}{c}}\toprule
$k$ & 0&1&2&3&4&5&6&7&8&9\\
$\dim R_k$ & 1&6&21&56&\textbf{120}&216&336&456&546&\textbf{580}\\\midrule
$k$ & 10&11&12&13&14&15&16&17&18&\\
$\dim R_k$ & 546&456&336&216&\textbf{120}&56&21&6&1&\\\bottomrule
\end{tabular}
\caption{Coefficients of $(1+t+t^2+t^3)^6$. Sum $=4096$. Bold: the Griffiths degrees.}\label{tab:series}
\end{table}
For example $R_4$ is spanned by degree-4 monomials in six variables with all exponents $\le 3$:
$\binom{9}{5}-6=126-6=120$ (only the six $x_i^4$ are excluded).
\section{Griffiths residues and the Hodge numbers}\label{sec:hodge}
\begin{theorem}[Griffiths~\cite{Griffiths69}, cited]
For a smooth hypersurface $X=V(G)\subset\PP^{n+1}$ of degree $d$, $n=\dim X$,
\[ H^{n-q,q}_{\mathrm{prim}}(X)\;\cong\;R_{d(q+1)-(n+2)}. \]
\end{theorem}
Here $n=4$, $d=5$, so $h^{4-q,q}_{\mathrm{prim}}=\dim R_{5q-1}$:
\begin{center}\begin{tabular}{cccc}\toprule
$q$ & degree $5(q+1)-6$ & $h^{4-q,q}_{\mathrm{prim}}$ & Lean\\\midrule
0 & $-1$ & 0 & \texttt{quintic\_h40\_zero}\\
1 & 4 & 120 & \texttt{quintic\_primitive\_hodge}\\
2 & 9 & 580 & \texttt{quintic\_primitive\_hodge}\\
3 & 14 & 120 & \texttt{quintic\_primitive\_hodge}\\
4 & 19 & 0 & (beyond socle degree 18)\\\bottomrule
\end{tabular}\end{center}
The total $0+120+580+120+0=820$ agrees with the Euler-characteristic computation of
Section~\ref{sec:setup}, an independent check. Adding $h^2$ gives the middle row of the Hodge diamond
\[ (h^{4,0},h^{3,1},h^{2,2},h^{1,3},h^{0,4})=(0,120,581,120,0),\qquad b_4=821. \]
\begin{remark}[What is and is not formalised]
Lean verifies the \emph{counting}: the coefficient list of $(1+t+t^2+t^3)^6$ and its entries at the
Griffiths degrees. The identification of $\dim R_k$ with a Hodge number is Griffiths' theorem, which is
cited. In particular $h^{3,1}=120\neq0$; only $h^{4,0}$ (and $h^{0,4}$) vanish.
\end{remark}
\section{Computing the Hilbert series in Lean}\label{sec:conv}
Lean represents a polynomial in $t$ by its coefficient list. The definitions (file \texttt{Hodge/FermatCubicHost.lean}) are
\begin{verbatim}
def conv (p q : List Nat) : List Nat :=
  (List.range (p.length + q.length - 1)).map fun k =>
    ((List.range (k + 1)).map fun i => p.getD i 0 * q.getD (k - i) 0).sum
def jacobianSeries (d : Nat) : List Nat :=
  (List.range 6).foldl (fun acc _ => conv acc (List.replicate (d - 1) 1)) [1]
def primHodge (d q : Nat) : Nat := (jacobianSeries d).getD (d * (q + 1) - 6) 0
\end{verbatim}
\texttt{conv} is the Cauchy product, so \texttt{jacobianSeries d} is the coefficient list of
$(1+t+\dots+t^{d-2})^6$, the Hilbert series of $\CC[x_0,\dots,x_5]/(x_0^{d-1},\dots,x_5^{d-1})$, and
\texttt{primHodge d q} is its coefficient in Griffiths degree $d(q+1)-6$. For $d=5$ the kernel evaluates
these closed terms by \texttt{decide}: \texttt{quintic\_series} (length 19, sum 4096) and
\texttt{quintic\_primitive\_hodge} ($120,580,120$). The same functions give $(1,20,1)$ for the cubic fourfold and
$(426,1751,426)$ for the sextic fourfold (\texttt{cubic\_primitive\_hodge}, \texttt{sextic\_primitive\_hodge}),
which are the known values and serve as sanity checks.

\paragraph{What the list computation does not prove.} Lean does not formalise the isomorphism between the
quotient ring $K[x]/J$ and the tensor product, nor its Hilbert series as a statement about a Mathlib
graded algebra; the link is the elementary monomial-basis argument of Section~\ref{sec:jac}. For $\psi=2$ the
identification uses that the six partials of a smooth quintic form a regular sequence (cited), and is
confirmed by Singular.

\section{The standalone certificate repository}
\url{https://github.com/BenFrohman/FermatQuinticJacobian} contains \texttt{FermatQuinticJacobian.lean},
which defines $G\in\mathbb Z[x_0,\dots,x_5]$ as a Mathlib \texttt{MvPolynomial} and proves
\texttt{Hodge.FermatQuintic.partial\_off} ($\partial_jx_i^5=0$ for $j\neq i$),
\texttt{partial\_on} ($\partial_ix_i^5=5x_i^4$) and \texttt{partial\_eq} ($\partial_iG=5x_i^4$).
Its CI audits these 3 declarations: each depends on \texttt{[propext, Classical.choice, Quot.sound]}.

\section{Singular members have infinite Jacobian rings}
For $X_1$ the partials have the common zero $p=(1,1,1,1,1,0)$ (Lean), so $K[x]/J(G_1)$ is not Artinian and
there is no Griffiths description; this is why the Hodge numbers above are claimed only for smooth members.
For $X_2$, Lean's \texttt{DworkQuinticPsi2.smooth} says the partials have only the trivial common zero, so
$J(G_2)$ is generated by a regular sequence and $R(G_2)$ has the same Hilbert series as the Fermat case,
which Singular confirms degree by degree.

\section{What is proved, computed, and cited}
\begin{center}\small
\begin{tabular}{p{0.55\linewidth}l}\toprule
Statement & Status\\\midrule
$\partial_iG=5x_i^4$; partials vanish together only at 0 (char.\ 0) & Lean\\
Coefficient list of $(1+t+t^2+t^3)^6$: length 19, sum 4096, entries 120, 580, 120 & Lean (\texttt{decide})\\
$R(G)$ has that Hilbert series & elementary (monomial basis)\\
Same series for every smooth quintic & cited (regular sequence)\\
$h^{4-q,q}_{\rm prim}=\dim R_{5q-1}$ & cited (Griffiths)\\
$X_1$ singular; $X_2$ smooth & Lean\\
Hilbert function of $R(G_2)$ & Singular (exact over $\QQ$)\\
$X_2$ contains no plane; $\ge 240$ Hodge classes & Singular; bookkeeping + cited theory\\
Hodge conjecture for $X_2$ & \textbf{open}\\\bottomrule
\end{tabular}\end{center}
\section{The Dwork-type pencil}\label{sec:dwork}
Consider the pencil
\[ G_\psi=\sum_{i=0}^5 x_i^5-5\psi\,x_0x_1x_2x_3x_4,\qquad X_\psi=V(G_\psi)\subset\PP^5 .\]
The product omits $x_5$, so every member is homogeneous of degree 5 (an earlier proposal using
$x_0\cdots x_5$ is \emph{not} homogeneous and does not define a hypersurface in $\PP^5$; this is also
recorded in Lean as \texttt{FermatQuintic.G\_not\_euler\_degree5}). $X_0$ is the Fermat quintic.

\subsection{$\psi=1$ is singular}
\begin{proposition}[Lean: \texttt{DeformedQuinticHost.singular\_at\_p}, \texttt{not\_smooth}]
The point $p=[1:1:1:1:1:0]$ lies on $X_1$ and all partial derivatives of $G_1$ vanish at $p$.
Hence $X_1$ is singular.
\end{proposition}
\begin{proof}
$G_1(p)=5-5=0$; $\partial_iG_1=5x_i^4-5\prod_{j\le4,j\neq i}x_j$ for $i\le4$ equals $5-5=0$ at $p$; and $\partial_5G_1=5x_5^4=0$.
\end{proof}
More generally $X_\psi$ is singular when $\psi^5=1$ (the points $[\zeta^{a_0}:\dots:\zeta^{a_4}:0]$);
only $\psi=1$ is formalised. The Griffiths machinery of Section~\ref{sec:hodge} does not apply to $X_1$.

\subsection{$\psi=2$ is smooth}
\begin{theorem}[Lean: \texttt{DworkQuinticPsi2.smooth}]\label{thm:psi2}
Over a field $K$ of characteristic $0$, the partials of $G_2=\sum x_i^5-10\,x_0\cdots x_4$ vanish
simultaneously only at $0$. Hence $X_2$ is a smooth quintic fourfold.
\end{theorem}
\begin{proof}
The partials are $5x_5^4$ and $5x_i^4-10\prod_{j\le 4,j\ne i}x_j$ ($i\le4$). So $x_5=0$ and
$x_i^4=2\prod_{j\ne i}x_j$. Put $P=x_0x_1x_2x_3x_4$. Multiplying the five relations,
$P^4=\prod x_i^4=2^5\prod_i\prod_{j\ne i}x_j=32P^4$, so $31P^4=0$ and $P=0$. Then
$x_i^5=x_i\cdot 2\prod_{j\ne i}x_j=2P=0$ for each $i\le4$. This is exactly the Lean proof
(it uses \texttt{linear\_combination} for each step).
\end{proof}
By Theorem~\ref{thm:psi2}, Section~\ref{sec:hodge} applies verbatim to $X_2$: primitive Hodge numbers
$0,120,580,120,0$ (Lean arithmetic: \texttt{dwork\_primitive\_hodge}). Independently, Singular computed
the Hilbert function of $\QQ[x]/J(G_2)$ as $1,6,21,56,120,216,336,456,546,580,546,\dots,1$ (total
4096), matching Table~\ref{tab:series} (script \texttt{scripts/dwork\_quintic\_psi2/hilb.sing}).

\subsection{Computations outside Lean for $\psi=2$}
\begin{itemize}
\item \textbf{No planes.} A plane in $\PP^5$ is the row space of a $3\times6$ matrix in reduced form.
$S_5$ (permuting $x_0,\dots,x_4$) preserves $G_2$, so the 20 Grassmannian charts reduce to two orbit
representatives (pivots $\{0,1,2\}$ and $\{0,1,5\}$). Restricting $G_2$ gives 21 equations in 9
unknowns; a Gr\"obner basis over $\QQ$ is $\{1\}$, so $X_2$ contains no plane over $\CC$. The same code
finds the planes of the Fermat quintic (sanity check). (Script \texttt{planes.sing}; exact, not Lean.)
\item \textbf{Hodge classes.} Let $H=\{a\in(\mathbb Z/5)^6: a_0+\dots+a_4\equiv0\}/\text{diagonal}$ (order 625)
act by $x_i\mapsto\zeta^{a_i}x_i$; it preserves every $G_\psi$. Deforming to $\psi=0$, the 820 Fermat
characters group into 500 $H$-eigenspaces (240 of rank 1, 200 of rank 2, 60 of rank 3). The 240 rank-1
eigenspaces are, with all Galois conjugates, of type $(2,2)$, so they give \emph{at least 240 primitive
rational Hodge classes on $X_\psi$ for every smooth $\psi$}, including $\psi=2$ (script \texttt{chars.py};
bookkeeping plus cited theory). The exact rank at $\psi=2$ is not computed.
\end{itemize}

\subsection{Status of the Hodge conjecture for $X_2$: open}
For the Fermat quintic the Hodge conjecture is known (Shioda~\cite{Shioda79}, Ran~\cite{Ran80}: prime degree),
with cycles built from linear subspaces. At $\psi=2$ those planes do not deform ($X_2$ contains none), and
$G_2$ has 7 monomials in 6 variables, so it is not of Delsarte type. Showing that the 240 classes above stay
algebraic along the pencil is an instance of the \emph{variational Hodge conjecture}, which is open. We found no
published result for $\psi=2$. \textbf{No claim of the Hodge conjecture for $X_2$ is made.}
\section{Lean formalisation}\label{sec:lean}
\subsection{Where the proofs live}
All machine-checked statements cited in this paper are in the repository
\url{https://github.com/BenFrohman/HODGE} (Lean~4 \texttt{v4.22.0}, Mathlib \texttt{v4.22.0}), in three files:
\begin{itemize}
\item \texttt{Hodge/QuinticHost.lean} (Fermat quintic, merged in PR~\#26);
\item \texttt{Hodge/DeformedQuinticHost.lean} (Dwork member $\psi=1$, PR~\#26);
\item \texttt{Hodge/DworkQuinticPsi2.lean} (Dwork member $\psi=2$, PR~\#28).
\end{itemize}
The counting functions \texttt{conv}, \texttt{jacobianSeries} and \texttt{primHodge} are defined in
\texttt{Hodge/FermatCubicHost.lean} and reused. The polynomial is an honest
\texttt{MvPolynomial (Fin 6) K} over an arbitrary field \texttt{K}; partial derivatives are Mathlib's
\texttt{MvPolynomial.pderiv}. None of the files uses \texttt{sorry}, \texttt{native\_decide} or a
user-declared \texttt{axiom}; the repository CI enforces this with a source scan.

\subsection{Theorem inventory}
\begin{center}\small
\begin{longtable}{p{0.42\linewidth}p{0.52\linewidth}}
\toprule
Lean declaration & Statement\\\midrule
\texttt{QuinticHost.G\_isHomogeneous} & $G=\sum_{i=0}^5 x_i^5$ is homogeneous of degree 5.\\
\texttt{QuinticHost.pderiv\_G} & $\partial G/\partial x_i = 5x_i^4$.\\
\texttt{QuinticHost.smooth} & If $\mathrm{char}\,K=0$ and all $\partial_iG(x)=0$ then $x=0$.\\
\texttt{QuinticHost.quintic\_series} & The coefficient list of $(1+t+t^2+t^3)^6$ has length 19 and sum 4096.\\
\texttt{QuinticHost.quintic\_primitive\_hodge} & $R_4=120$, $R_9=580$, $R_{14}=120$.\\
\texttt{QuinticHost.quintic\_h40\_zero} & $5\cdot 1<6$ (so the Griffiths degree for $h^{4,0}$ is negative).\\
\texttt{QuinticHost.quintic\_h31\_ne\_zero} & $R_4\neq 0$.\\
\texttt{DeformedQuinticHost.G\_isHomogeneous} & $G_1=\sum x_i^5-5x_0x_1x_2x_3x_4$ is homogeneous of degree 5.\\
\texttt{DeformedQuinticHost.singular\_at\_p} & $G_1(p)=0$ and $\partial_iG_1(p)=0$ for all $i$, $p=(1,1,1,1,1,0)$.\\
\texttt{DeformedQuinticHost.not\_smooth} & The Jacobian criterion fails for $G_1$.\\
\texttt{DworkQuinticPsi2.G\_isHomogeneous} & $G_2=\sum x_i^5-10x_0x_1x_2x_3x_4$ is homogeneous of degree 5.\\
\texttt{DworkQuinticPsi2.eval\_pderiv} & Explicit formulas for the six partials of $G_2$.\\
\texttt{DworkQuinticPsi2.smooth} & If $\mathrm{char}\,K=0$ and all $\partial_iG_2(x)=0$ then $x=0$.\\
\texttt{DworkQuinticPsi2.psi\_not\_root} & $2^5\neq 1$.\\
\texttt{DworkQuinticPsi2.dwork\_primitive\_hodge} & $R_4=120$, $R_9=580$, $R_{14}=120$ for the series $(1+t+t^2+t^3)^6$.\\
\bottomrule
\end{longtable}\end{center}

\subsection{Axiom audit}
The HODGE CI runs \texttt{lake env lean AxiomAudit.lean} (one \texttt{\#print axioms} per declaration)
and fails unless every declaration depends only on Lean's standard axioms
\texttt{propext}, \texttt{Classical.choice}, \texttt{Quot.sound}, or on none. Output on \texttt{main} (the $\psi=2$ declarations were added to the audit in HODGE PR~\#31):
\begin{center}\small
\begin{tabular}{p{0.52\linewidth}p{0.40\linewidth}}\toprule
Declarations & \texttt{\#print axioms}\\\midrule
\texttt{QuinticHost.\{G\_isHomogeneous, pderiv\_G, smooth\}} & \texttt{[propext, Classical.choice, Quot.sound]}\\
\texttt{QuinticHost.\{quintic\_series, quintic\_primitive\_hodge, quintic\_h31\_ne\_zero\}} & \texttt{[propext]}\\
\texttt{QuinticHost.quintic\_h40\_zero} & none\\
\texttt{DeformedQuinticHost.*} (3 declarations) & \texttt{[propext, Classical.choice, Quot.sound]}\\
\texttt{DworkQuinticPsi2.\{G\_isHomogeneous, eval\_pderiv, smooth\}} & \texttt{[propext, Classical.choice, Quot.sound]}\\
\texttt{DworkQuinticPsi2.\{psi\_not\_root, dwork\_primitive\_hodge\}} & \texttt{[propext]}\\
\bottomrule\end{tabular}\end{center}
\texttt{Classical.choice} enters through Mathlib's polynomial and field library, not through any
step specific to this paper. The counting theorems are closed by \texttt{decide} and use only \texttt{propext}.
\section{Reproducibility}
\begin{itemize}
\item Build: \texttt{git clone https://github.com/BenFrohman/HODGE; cd HODGE; lake exe cache get; lake build}.
\item Axiom audit: \texttt{lake env lean AxiomAudit.lean}; CI workflow \texttt{Lean certificate}.
\item Singular scripts for $\psi=2$: \texttt{scripts/dwork\_quintic\_psi2/} (\texttt{hilb.sing}, \texttt{planes.sing},
\texttt{fermat.sing}, \texttt{chars.py}) with outputs committed next to them.
\item Standalone certificate repository: \url{https://github.com/BenFrohman/FermatQuinticJacobian}, which carries its own
CI with the same axiom check.
\end{itemize}

\section*{Scope}
This paper proves smoothness statements and Jacobian-ring counts in Lean and cites Griffiths' theorem for
the Hodge-theoretic interpretation. It makes no claim about algebraic cycles on the quintic fourfold beyond
the cited Fermat results, and it does not claim the Hodge conjecture for any non-Fermat member of the pencil.

\begin{thebibliography}{9}
\bibitem{Griffiths69} P.~A. Griffiths, On the periods of certain rational integrals I, II, \emph{Ann. of Math.} 90 (1969), 460--541.
\bibitem{Shioda79} T.~Shioda, The Hodge conjecture for Fermat varieties, \emph{Math. Ann.} 245 (1979), 175--184.
\bibitem{Ran80} Z.~Ran, Cycles on Fermat hypersurfaces, \emph{Compositio Math.} 42 (1980), 121--142.
\bibitem{Voisin} C.~Voisin, \emph{Hodge Theory and Complex Algebraic Geometry I, II}, Cambridge Univ. Press, 2002--2003.
\bibitem{Mathlib} The mathlib Community, The Lean mathematical library, \emph{CPP 2020}, 367--381.
\end{thebibliography}
\end{document}
